\section{立憲ガバナンスと共進化(ARI-RQGM)}\label{sec:rqgm}

ここまでの各章は、ARI の\term{制度}(institution)——研究の方向を
提案し、ノードをレビューし、品質を裁定する、プロンプトで定義された
役割群——を、ランの間じゅう固定のものとして扱ってきた。本章で述べる
\term{ARI-RQGM} は、その制度を、自らは決して進化しない憲法の下で
進化させる、オプトイン式の実行モードである。同モードは
\cref{sec:exploration} の最良優先木探索を変更しないまま、エポック
ベースのガバナンス層(\texttt{ari.rqgm} パッケージ)で包む:
プロンプトとコンポーネントは競争し、攻撃し、防御し、採用され、
あるいは退役する——しかし制度のあらゆる状態変更は、コードに凍結
された規則テーブルに照らして、決定的で LLM を用いないカーネルに
より検証される。既定のモードは依然として統治なしのループであり、
ガバナンスを未設定のままにすればガバナンスオブジェクトは一切構築
されず、チェックポイントは先行章の述べるものとバイト単位で同一で
ある。オプトインはラン単位の決定であり、ラン開始前になされ、
チェックポイントに記録される。

\subsection{動機と系譜}\label{sec:rqgm:motivation}

制度的役割はそれ自体プロンプトで定義された成果物であるから、それら
を進化させれば、探索ループが実験から引き出すのと同種の改善が見込
める。しかし、制約のない自己改変こそ G\"odel マシン系の文献が警告
するものである:自身の評価器を書き換えられるコンポーネントは、
いずれ自分自身に報酬を与える。Red Queen G\"odel Machine
\cite{iacob2026red} は共進化という対抗錘を加える——エージェントが
改善するのは、敵対者と評価器がそれに対して共進化する\emph{から
こそ}である——ものであり、ARI-RQGM の名と中心的力学はこの系統の
研究に由来する。だが共進化だけでは、研究システムにとって重要な
故障モード——スコアの水増し、メトリックゲーミング、証拠の捏造、
役割間の共謀——を抑え込むことはできない。

そこで ARI-RQGM は、共進化ループを、決して進化しない固定の憲法層
の下に置く。全体を貫く設計姿勢は\emph{制度はブロックしても、研究は
決してブロックしない}である:ガバナンスは自身の状態変更——採用、
制裁、レジストリ書き込み——を拒否できるが、ノードの実行を拒否する
ことは決してなく、探索経路上のすべてのガバナンスフックは fail-open
であるため、ガバナンスの故障は実験ではなくガバナンスを劣化させる。
これは \cref{sec:overview:principles} の決定性原則を制度の変更に
まで拡張したものである:カーネルの判定、遷移の解決、フロンティア
修復は純関数であり——乱数も壁時計による判定もない——統治なしの
ランが同じ木へ再生されるのとまったく同様に、統治されたランは同じ
制度史へと再生される。

\subsection{三つの層、四つのファサード}\label{sec:rqgm:layers}

ガバナンスはシステムを三つの層へ成層化する。層を降りるにつれ、
物事を変更する権限は厳密に減少し、変更される自由は厳密に増加する。

\begin{description}
  \item[\textbf{Layer 0(憲法層、固定)。}] 決定的カーネル、固定の
        結果検証経路、そして append-only の監査ログ。これらは決して
        進化しない;それらが執行する規則テーブルはコードに存在し、
        憲法ハッシュ(\cref{sec:rqgm:constitution})でピン留めされ、
        来歴のためだけにコンポーネントレジストリへ登録される。
  \item[\textbf{Layer 1(制度層)。}] 進化可能なサービス役割:
        ジェネレータ、レビュアー、敵対者、弁護人、ジャッジ、
        ルータ。これらの役割のプロンプトは、\cref{sec:rqgm:prompts}
        のライフサイクルを通じて、エポック境界においてのみ置き換え
        られうる。
  \item[\textbf{Layer 2(メタ層)。}] Layer~1 を進化させる
        エージェント——プロンプト変異器とクリーンルーム生成器、
        およびより狭い補助エージェント——であり、それ自身も統治
        され、変更対象の層より厳密に小さい権限しか持たない:メタ
        エージェントは置き換えを提案できるが、任命は決してできない。
\end{description}

この機構は四つのファサードが所掌する。\term{憲法カーネル}
(constitutional kernel)は、レコード、ハッシュ、ケイパビリティ、
遷移、役割分離、消去、監査ログ完全性、クリーンルーム入力、
コンテキストスコープを検証する——12 の閉じたエントリポイントで
あり、どれもモデルを呼ばず、時計を読まない。\term{ガバナンス
オーケストレータ}(governance orchestrator)はエポック境界の監査
(\cref{sec:rqgm:boundary})を実行して助言的なレポートを生成する;
レジストリに書くことは決してない。\term{レジストリ遷移エンジン}
(registry transition engine)はコンポーネントとプロンプトの
ステータスの唯一の書き手であり、固定の遷移テーブルに照らして変更を
解決する。\term{フロンティア修復エンジン}(frontier repair engine)
は、退役の後に探索フロンティアを復元する(\cref{sec:rqgm:boundary})。
その周囲にノード単位のサービス群が位置する:提案ルータ、敵対
ラウンド、プロンプト進化パイプライン、クリーンルームコーディネータ、
そして、ノードの実行には決して触れずにあらゆるガバナンス決定点を
縮退させうる予算マネージャである。

\subsection{起動:モード解決と創設登録}\label{sec:rqgm:boot}

\begin{figure}[t]
  \centering
  \resizebox{\columnwidth}{!}{% F22: full lifecycle of a governed (ari_rqgm) run, set in the ARI house
% design language (_style.tex). Three phase columns, F21-style: BOOT
% (lavender / ideation fill) — mode resolution, the one-time founding
% registration transaction (25 prompts + 12 components), and the first
% epoch freeze; GOVERNED SEARCH (peach / execution fill) — the per-node
% hooks (proposal round, ungoverned node execution, governance-level
% assignment, adversarial round) inside a dashed loop enclosure; EPOCH
% BOUNDARY (pale green / gate fill) — audit, kernel-validated transition,
% frontier repair + clean room, and the next freeze. The dotted margin
% path at the bottom is the reflux into the next epoch's search. No
% accent edge: this figure has no single emphasis role (F23 carries it).
%
% style.tikzstyles is loaded by the preamble / the standalone wrapper;
% the ari* house styles come from _style.tex. The build cwd differs per
% entry point, so resolve _style.tex from each known prefix.
\InputIfFileExists{../shared/figures/tikz/_style.tex}{}{%
  \InputIfFileExists{shared/figures/tikz/_style.tex}{}{%
    \InputIfFileExists{report/shared/figures/tikz/_style.tex}{}{%
      \InputIfFileExists{_style.tex}{}{%
        \errmessage{F22: cannot locate figures/tikz/_style.tex}}}}}
\begin{tikzpicture}[aricanvas]

  % --- column 1: boot (mode + founding registration) -------------------
  \node[ariboxwide, fill=phaseidea]           (mode)     {\figlabel{F22.mode}};
  \node[ariboxwide, fill=phaseidea, below=of mode] (founding) {\figlabel{F22.founding}};
  \node[ariboxwide, fill=phaseidea, below=of founding] (freezeo) {\figlabel{F22.freeze0}};
  \node[ariphase, above=0.4cm of mode]        (h1)       {\figlabel{F22.head.boot}};
  \pic at ([xshift=-2.5mm]h1.west) {ariglyph bulb};
  \draw[ariarrow] (mode)     -- (founding);
  \draw[ariarrow] (founding) -- (freezeo);

  % --- column 2: governed search (per-node hooks, loop-enclosed) -------
  \node[aribox exp, right=of mode]            (proposal) {\figlabel{F22.proposal}};
  \node[aribox exp, below=of proposal]        (execute)  {\figlabel{F22.execute}};
  \node[aribox exp, below=of execute]         (level)    {\figlabel{F22.level}};
  \node[aribox exp, below=of level]           (advers)   {\figlabel{F22.adversarial}};
  \node[ariloop, fit=(proposal)(execute)(level)(advers)] (loopgov) {};
  \node[ariphase, above=0.65cm of proposal]   (h2)       {\figlabel{F22.head.search}};
  \pic at ([xshift=-2.5mm]h2.west) {ariglyph flask};
  \node[arinote, anchor=south west] at ([xshift=6mm]loopgov.north) (ttl) {\figlabel{F22.loop}};
  \pic at ([xshift=-2.5mm]ttl.west) {ariglyph loop};
  \draw[ariarrow] (proposal) -- (execute);
  \draw[ariarrow] (execute)  -- (level);
  \draw[ariarrow] (level)    -- (advers);

  % --- column 3: epoch boundary ----------------------------------------
  \node[aribox gate, right=of proposal]       (audit)    {\figlabel{F22.audit}};
  \node[aribox gate, below=of audit]          (transit)  {\figlabel{F22.transition}};
  \node[aribox gate, below=of transit]        (repair)   {\figlabel{F22.repair}};
  \node[aribox idea, below=of repair]         (freezen)  {\figlabel{F22.freeze1}};
  \node[ariphase, above=0.4cm of audit]       (h3)       {\figlabel{F22.head.boundary}};
  \pic at ([xshift=-2.5mm]h3.west) {ariglyph gate};
  \draw[ariarrow] (audit)   -- (transit);
  \draw[ariarrow] (transit) -- (repair);
  \draw[ariarrow] (repair)  -- (freezen);

  % --- handoffs between the phase columns -------------------------------
  \draw[ariarrow, rounded corners=5pt]
    (freezeo.east) -- ([xshift=8mm]freezeo.east) |- (proposal.west);
  \draw[ariarrow, rounded corners=5pt]
    (advers.east) -- ([xshift=8mm]advers.east)
    |- node[arinote, pos=0.3] {\figlabel{F22.note.trigger}} (audit.west);

  % --- reflux: the new epoch keeps searching ----------------------------
  \draw[arifeedback]
    (freezen.south) -- ([yshift=-7mm]freezen.south) -| (loopgov.south);
  \node[arinote] at ([yshift=-7mm]loopgov.south east) {\figlabel{F22.note.next}};

\end{tikzpicture}
}
  \caption{統治されたランのライフサイクル。起動時に実行モードを
    解決し、一度限りの創設登録トランザクション(プロンプト 25 件、
    コンポーネント 12 件)を実行し、創設アクティブ集合をエポック~000
    として凍結する。エポックの間、ノード単位のフックが提案を記録し、
    ガバナンスレベル(固定チェックのみの L0 から完全裁定付きの L3
    まで)を割り当て、イベント駆動の敵対ラウンドを実行する;ノード
    の実行そのものは決して拒否されない。エポックのノード数トリガが
    発火すると、境界は助言的監査、カーネル検証付きの遷移、そして
    クリーンルーム再生成を伴うフロンティア修復を実行し、次エポック
    のアクティブ集合を凍結して、探索は継続する。}
  \label{fig:rqgm-lifecycle}
\end{figure}

統治されたランは三段階で起動する(\cref{fig:rqgm-lifecycle})。
第一に\term{モード解決}(mode resolution):モードは二鍵
インターロックであり——モードセレクタと冗長な有効化フラグが一致
しなければならず、環境変数による上書きは代入前に検証される——
不一致があれば、推測する代わりに警告の上で統治なしの既定へ
フォールバックする。解決されたモードとその来歴(設定、環境、
または再開)はチェックポイントに記録される。

第二に、ランタイムは探索戦略を純粋委譲のシムで包む:選択、拡張、
枝刈りは \cref{sec:exploration} とまったく同様に振る舞い、
ガバナンスはその周囲にベストエフォートのフックとして付随する。

第三に——新規の統治チェックポイントに限り——\term{創設登録}
(founding registration)トランザクションが制度を初期化する。
コードに凍結された二つのテーブルから、一つの prepare イベントに
続いて 25 件のプロンプト登録と 12 件のコンポーネント登録が行われ、
最後に単一のコミットが置かれる。創設コンポーネント 12 件とは、
七つの敵対者バリアント(攻撃タイプごとに一つ、すべて敵対者役割を
共有する)、成果物ジャッジ、弁護人、提案ルータ、そして二つのメタ層
エージェント(プロンプト変異器とクリーンルーム生成器)である;
創設プロンプト 25 件はこれに加えて、ガバナンスアクター(監査人、
ガバナンス弁護人、ガバナンス裁定者)、レビュアーテンプレート、
\cref{sec:exploration} の探索テンプレートを覆い、各プロンプトは
そのテンプレートバイト列のハッシュとともに登録される。
トランザクションは次いで、創設アクティブ集合——アクティブな
コンポーネント、アクティブなプロンプトハッシュ、効用ポリシー——を
エポック~000 の不変状態として凍結する。そのフィンガープリントは
意図的にタイムスタンプを除外しており、同一の制度は同一の
フィンガープリントを持つ。創設は write-once である:再開時には
イベントログの再生によってレジストリが再構築され、コミット済みの
登録が再実行されることはなく、コミット前に中断された創設
トランザクションは不可視であって、単に再実行される。

\subsection{ノード単位のガバナンス}\label{sec:rqgm:node}

エポックの間、探索は \cref{sec:exploration} とまったく同様に走る;
ガバナンスは三つのベストエフォートフックを通じてそれを観測する。

\paragraph{提案。}
拡張ステップが候補方向を生成すると、提案ルータは、「すべてを保存
する」原則に基づき、生成コンテキストの全体からなる完全な\term{提案
レコード}(proposal record)をアーカイブする。一方、探索自体が
受け取るのは上限付きの\term{提案要約ビュー}(proposal summary
view)だけである。要約は探索が見る\emph{唯一の}表面であるため、
提案ジェネレータは、探索ループへのインタフェースを広げることなく
豊かになることができ、アーカイブされたレコードは監査と再生のために
利用可能であり続ける。

\paragraph{ガバナンスレベル。}
評価された各ノードには、純粋なトリガのはしごによって、L0(固定
チェックのみ)から L3(完全な裁定付き)までのガバナンスレベルが
割り当てられる。根拠は決定的なシグナル——フロンティア上位スコアへ
の所属、親に対する異常なスコアジャンプ、論文候補フラグ、新規性の
主張、評価器の低い確信度——である。レベルが決めるのはノードが
受けるガバナンスの注意の量であって、ノードが実行されるか否かでは
決してない。すべての割り当ては監査ログへ追記され、エポックごとの
ガバナンス予算が枯渇しかけたときには、予算マネージャが実効レベルに
上限を課すことがある——縮退するのはガバナンスであって、ノードでは
ない。

\paragraph{敵対ラウンド。}
同じトリガ族に該当したノードは\term{敵対ラウンド}(adversarial
round)も受ける:七つの敵対者タイプ——過大主張、証拠ギャップ、
メトリックゲーミング、再現性、先行研究、コスト爆発、プロンプト
インジェクション——のいずれかが、ノードの\emph{成果物}を攻撃する。
敵対者がコンポーネントを攻撃することは決してない;攻撃レコードには、
構成からしてコンポーネントのフィールドが存在しない。生の攻撃は
スコアに対して不活性である:弁護人が応答し、成果物ジャッジが攻撃を
認容して初めて\term{検証済み攻撃レコード}(validated attack
record)が存在し、検証済みの攻撃だけが、エポック凍結されたペナルティ
ポリシーの下で有界の効用ペナルティへ算入される。ペナルティ前の
スコアは加法的なメトリックキーに保存される。ジャッジの故障は既定で
\emph{無効}——すなわちペナルティなし——に倒れる。検証済みの攻撃は
また、後に候補プロンプトを試験するコーパスである\term{敵対的
リプレイプール}(adversarial replay pool)へも入る
(\cref{sec:rqgm:prompts})——これが Red Queen 的な結合である:
今日の検証済み攻撃は、明日の入学試験である。ラウンドはノードごとに
冪等である——ケースログ内のマーカーにより、再開されたランが処理
済みのノードを再攻撃することは決してない。

\subsection{エポック境界}\label{sec:rqgm:boundary}

\begin{figure}[t]
  \centering
  \resizebox{\columnwidth}{!}{% F23: the epoch-boundary governance pipeline, set in the ARI house design
% language (_style.tex). Left column: the nine-step advisory audit
% (observe .. self-audit, numbered badges; the dashed loop enclosure marks
% it read-only w.r.t. registries). Right column: the report, the
% transition engine's resolution, the constitutional kernel's validation
% of every edge against the fixed T1-T21 table, the four-event boundary
% transaction, and frontier repair. The single accent edge (the rationed
% emphasis role of this figure) is the kernel's fail-closed block of an
% illegal transition.
%
% style.tikzstyles is loaded by the preamble / the standalone wrapper;
% the ari* house styles come from _style.tex. The build cwd differs per
% entry point, so resolve _style.tex from each known prefix.
\InputIfFileExists{../shared/figures/tikz/_style.tex}{}{%
  \InputIfFileExists{shared/figures/tikz/_style.tex}{}{%
    \InputIfFileExists{report/shared/figures/tikz/_style.tex}{}{%
      \InputIfFileExists{_style.tex}{}{%
        \errmessage{F23: cannot locate figures/tikz/_style.tex}}}}}
\begin{tikzpicture}[aricanvas]

  % --- column 1: the nine-step epoch audit (advisory) -------------------
  \node[aribox exp]                        (observe)  {\figlabel{F23.observe}};
  \node[aribox exp, below=of observe]      (reliab)   {\figlabel{F23.reliability}};
  \node[aribox exp, below=of reliab]       (evidence) {\figlabel{F23.evidence}};
  \node[aribox exp, below=of evidence]     (prosec)   {\figlabel{F23.prosecute}};
  \node[aribox exp, below=of prosec]       (defense)  {\figlabel{F23.defense}};
  \node[aribox exp, below=of defense]      (adjud)    {\figlabel{F23.adjudicate}};
  \node[aribox exp, below=of adjud]        (pool)     {\figlabel{F23.replaypool}};
  \node[aribox exp, below=of pool]         (selfaud)  {\figlabel{F23.selfaudit}};
  \node[ariloop, fit=(observe)(reliab)(evidence)(prosec)(defense)(adjud)(pool)(selfaud)] (loopaud) {};
  \node[ariphase, above=0.65cm of observe] (h1)       {\figlabel{F23.head.audit}};
  \pic at ([xshift=-2.5mm]h1.west) {ariglyph flask};
  \node[arinote, anchor=north] at ([yshift=-2.5mm]loopaud.south) (ttl) {\figlabel{F23.note.advisory}};
  \draw[ariarrow] (observe)  -- (reliab);
  \draw[ariarrow] (reliab)   -- (evidence);
  \draw[ariarrow] (evidence) -- (prosec);
  \draw[ariarrow] (prosec)   -- (defense);
  \draw[ariarrow] (defense)  -- (adjud);
  \draw[ariarrow] (adjud)    -- (pool);
  \draw[ariarrow] (pool)     -- (selfaud);

  % Step badges 1..8 (the report itself is step 9).
  \node[aribadge] at (observe.north west)  {1};
  \node[aribadge] at (reliab.north west)   {2};
  \node[aribadge] at (evidence.north west) {3};
  \node[aribadge] at (prosec.north west)   {4};
  \node[aribadge] at (defense.north west)  {5};
  \node[aribadge] at (adjud.north west)    {6};
  \node[aribadge] at (pool.north west)     {7};
  \node[aribadge] at (selfaud.north west)  {8};

  % --- column 2: report -> transition -> kernel -> commit -> repair -----
  \node[aribox paper, right=2.4cm of observe] (report) {\figlabel{F23.report}};
  \node[aribox idea,  below=of report]        (resolve){\figlabel{F23.resolve}};
  \node[aribox gate,  below=of resolve]       (kernel) {\figlabel{F23.kernel}};
  \node[ariboxwide, fill=phasepaper, below=of kernel] (txn) {\figlabel{F23.txn}};
  \node[aribox exp,   below=of txn]           (repair) {\figlabel{F23.repair}};
  \node[arinote, below=0.2cm of repair]       (retn)   {\figlabel{F23.note.retire}};
  \node[ariphase, above=0.4cm of report]      (h2)     {\figlabel{F23.head.transition}};
  \pic at ([xshift=-2.5mm]h2.west) {ariglyph gate};
  \node[aribadge] at (report.north west) {9};
  \draw[ariarrow] (report)  -- (resolve);
  \draw[ariarrow] (resolve) -- (kernel);
  \draw[ariarrow] (kernel)  -- (txn);
  \draw[ariarrow] (txn)     -- (repair);

  % --- the audit hands its report across the column gap -----------------
  \draw[ariarrow, rounded corners=5pt]
    (selfaud.east) -- ([xshift=10mm]selfaud.east) |- (report.west);

  % --- the one accent role: an illegal edge is blocked fail-closed ------
  \node[ariboxwide, right=2.2cm of kernel] (blocked) {\figlabel{F23.blocked}};
  \draw[ariarrow accent] (kernel) -- (blocked);

\end{tikzpicture}
}
  \caption{エポック境界パイプライン。九段の監査は助言的であり、
    レジストリに対しては読み取り専用である;そのレポートは
    レジストリ遷移エンジンへ供給され、解決された遷移は、四イベント
    の境界トランザクションがコミットする前に、憲法カーネルによって
    固定の T1--T19 テーブルに照らしエッジ単位で検証される。不正な
    エッジは fail-closed でブロックされる——現職のアクティブ集合が
    引き続き奉仕し、ランは継続する。フロンティア修復が走るのは、
    コミットされた遷移がコンポーネントまたはプロンプトを退役させた
    場合に限られる。}
  \label{fig:rqgm-boundary}
\end{figure}

エポックは制度を凍結する;境界は、制度が変わりうる唯一の窓である。
境界トリガはノード数(既定ではエポックあたり新規 10 ノード)であり、
外側の探索反復の先頭でメインスレッド上で——ノードが飛行中のときには
決して——検査され、さらにラン終了時にもう一度検査される。これに
より、最終ノードが最後の反復で到着したエポックもトランザクショナル
に閉じ、トリガ未満の末尾エポックは単に開いたままとなる。ラン終了時
の検査が存在するのは、実際のランがその欠落を露呈したからである:
ちょうど最後のノードで到達した境界は、さもなければ決して発火せず、
そのエポックの監査は書かれないままになる。

\paragraph{監査。}
閉じようとするエポックはまず監査される(\cref{fig:rqgm-boundary})。
監査は九段のパイプラインであり、LLM を消費するすべての段に全面的な
決定的フォールバックがあるため、モデルなしの監査も完走する。
エポックのレコードから観測が収集され;信頼性評価が統治対象の各
コンポーネントを採点し;\term{証拠係}(evidence clerk)が——そして
係だけが——証拠バンドルを組み立てる。同一役割からの告発は憲法に
より証拠能力を持たず、除外される;続いて訴追があり、弾劾動議を提起
できるのは\term{監査人}(auditor)のみであって、エポックごとの訴追
予算に対して保証金を供託し、動議が認容されれば返還される——濫訴に
対する経済的ブレーキである;すべての動議に対して弁護が生成され;
裁定パネルが裁定を下し;リプレイプールが更新され;\term{ガバナンス
自己監査}(governance self-audit)が、ガバナンスアクター自身の
レコードをカーネルを通じて再検証する。結果は\term{ガバナンス
レポート}(governance report)——信頼性、動議、弁護、裁定、候補
評価、勧告、保証金の会計——であり、ハッシュ連鎖された監査ログへ
追記される。レポートは次段への助言的入力である;それを生成した
オーケストレータはレジストリに書くことができない。

\paragraph{遷移。}
レジストリ遷移エンジンは、レポートの勧告と候補評価を、一つの
エポック遷移——採用、制裁、退役——へと解決する。これは凍結された
エポック状態、レポート、レジストリ、設定されたしきい値の純関数で
あり、再実行すればバイト単位で同一である。コンポーネントと
プロンプトは、十のステータス(候補、検証済み、シャドー、試用
アクティブ、アクティブ、警告、保護観察、隔離、退役、追放)を、
十九の規則 T1--T19 からなる固定テーブルに沿って移動する。カーネルは
解決された遷移のすべてのエッジを検証する:テーブル外のエッジは
名前付きの違反 \texttt{CK-REG-001} であり、境界限定のエッジを緊急の
形式なしにエポック途中に刻印すれば \texttt{CK-REG-002}、遷移
エンジン以外が作成した遷移——たとえばジャッジによるレジストリ
書き込み——は \texttt{CK-ROL-901} である。ブロックされた遷移は
fail-closed で中止され、現職のアクティブ集合が引き続き奉仕する;
停止されるのはガバナンスであって、ランでは決してない。有効な遷移は、
四イベントの境界トランザクション——prepare、旧エポックの close、
新エポックの open、commit——を通じてコミットされる。したがって、
どの二イベント間でクラッシュしても、残るのは再開時に破棄され決定的
に再実行される接頭辞だけであり、同一遷移の重複コミットはガードされた
no-op である。境界限定の変更に対する唯一の例外は\term{緊急隔離}
(emergency quarantine;規則 T16)である:カーネル・クリティカルな
違反の現場で捕捉された奉仕中のコンポーネントはエポック途中でも隔離
でき、その緊急措置は他のあらゆる遷移と同様に記録され、検証される。

\paragraph{フロンティア修復。}
コミットされた遷移がコンポーネントまたはプロンプトを退役させた場合、
ランの過去には、いまや欠陥ありと判断された制度の下で生成された
レコードが含まれており、フロンティアがそれらに利息を払い続けては
ならない。修復エンジンは、退役したプロンプトハッシュの依存閉包——
そこから実質的に導出されたすべてのレコード——を計算し、影響を受けた
ノードを加法的なフラグ(stale、valid-for-frontier、消去理由)で
stale とマークし、ツリーファイルを通じて永続化する。消去は
\emph{論理のみ}である:物理的には何も削除されないため、履歴は監査
可能なまま残り、フロンティアのスコアリングだけがそれを無視する。
修復は役割を認識する:退役したレビュアーのスコアを単に\emph{消費
した}効用は、\emph{元の}エポックの凍結された重みの下で、生き残った
入力から再計算される。一方、退役した役割自身が生成したレコードは
端的に無効化される——そして元の重みが復元できないときは、ノードは
再スケールされるのではなく無効化されるため、いかなる数値も、異なる
ポリシーの下で黙って再導出されることはない。修復自体が検証に失敗
した場合、エンジンは保守的に縮退し、最終的には drain-only モード
——ランは保留中の作業を完了するが、それ以上は拡張しない——に至る。
クラッシュは決して起こさない。

\subsection{進化するスコア}\label{sec:rqgm:utility}

ここまでの各小節は、制度を進化させながら、一つのものを固定として
扱ってきた:\emph{スコア}そのもの——フロンティアを順位づける効用
ポリシーであり、軸の基底とその重み、合成規則、フロンティアスコア
リング規則、深さペナルティ、そして探索定数から成る。スコアをランの
間じゅう凍結するのは保守的な選択だが、Red Queen 系の研究
\cite{iacob2026red} が実際に行う選択ではない:その定義的な主張は、
各エポック境界でスコア全体が書き換えられる木探索である。統治された
効用進化はその隔たりを埋める。効用ポリシーはそれ自体が一個の進化
可能な役割となる——\cref{sec:rqgm:prompts} のプロンプト役割との
唯一の違いは、その現職がプロンプトテンプレートではなくポリシー文書
である点だ。効用ポリシーはエポック内では凍結され、エポックが走る間
フロンティアは一つの固定規則で順位づけられ、そして境界においてのみ、
プロンプトを採用するのと同じ遷移エンジンによって書き換えられうる。

専用の提案器が候補ポリシーを鋳造し、それ以外のことは何もしない。
その四つの既定変異は、境界のすでに抽象化された証拠に対する純粋な
算術であり——モデルも時計も乱数もない——ゆえに既定の書き換えは
完全に再現可能であり、しかもフロンティアの現在のスコアを読むことは
禁じられている:自らがすでに生み出したノードにおもねるよう調整
されたスコアは、まさに共進化的設計が防ぐために存在する自己言及的な
共謀である。提案器はレジストリ権限を持たず、それ自体が登録され制裁
されうるコンポーネントである——それが採点する役割の上にいないのと
同様に、スコアの上に絶対的な支配者はいない。

採用は遷移テーブルの一つの新しいエッジに乗る。検証とシャドー奉仕を
通過し、凍結されたリプレイ盤上で現職と少なくとも同等に採点される
後継ポリシーは、現職を\emph{追い越す}:規則~T20、唯一の active から
retired へのエッジであり、隔離を経由しない唯一の退役である。カーネル
はこれを効用ポリシー役割にのみ許す——あらゆる行動的コンポーネントは
保守的な制裁のみの置換モデルを保ち、それらにとって直接の退役は依然
として拒否される。追い越しは健全な現職を\emph{古い}ポリシーハッシュ
の下で退役させ、いまやすべてのノードが自らを採点したポリシーの
ハッシュを帯びているため、フロンティア修復(\cref{sec:rqgm:boundary})
は、たまたまペナルティを受けた部分集合だけでなく、退役したポリシーの
下で順位づけられたすべてのノードを無効化する。このエッジの追加は、
改正がそうあらねばならないように、憲法ハッシュをピン留めし直した
(\cref{sec:rqgm:constitution})。

一つの限界は、機構と同じくらい率直に述べるに値する。スコアの軸が
動的に導出され、固定の軸ごとの基底がピン留めされていないとき、
「少なくとも同等に採点される」というゲートには比較すべき安定な順序
が存在しないため、追い越しは、実証された改善ではなく、適法性と
非退化性へと縮退する。後継を\emph{厳密に}より良いと証明できる
ゲートは、採点の拠り所となるアンカー基底——採点対象とする真値——を
必要とする。それは \cref{sec:rqgm:paper} の論文パイプラインが供給
するが、研究品質に対する外部のオラクルを持たない探索フェーズは供給
しない。これが率直な Red Queen の姿勢である:真値がなければ、基準は
固定の丘を登るのではなく適法なポリシーの空間の内を転がるのであって、
我々はそれを実態以上のものとして提示しない。

\subsection{攻撃を説明責任あるコンポーネントへ束縛する}\label{sec:rqgm:targetbinding}

\cref{sec:rqgm:node} の敵対ラウンドは、裁定の後に検証済み攻撃
レコードを生成する。そのレコードは常に、自らが実装する\emph{役割}を
名指してきた——役割こそ攻撃が誠実に知りうる唯一のものである——が、
これまで説明責任ある\emph{コンポーネント}を名指すことはなく、ゆえに
検証済み攻撃から、信頼性カウントを経て、弾劾動議へと走るはずの連鎖
は、指し示す先を持たず不活性のままだった。理由は構造的である:七つ
の探索敵対者は、登録された現職を持たないノード生成役割が作成した
成果物を攻撃するため、原理的にすら説明責任を問うべきコンポーネント
が存在しなかった。

いまやジャッジが作成するレコードは、解決可能なときにはいつでも、
単一のレジストリ解決可能なターゲットを帯びる:攻撃が無効化する決定
を下した現職であり、レコードが鋳造される時点でエポック\emph{凍結}
されたアクティブ集合に照らして解決される——決してライブなレジストリ
読み出しではない。それは境界での採用の後には、先任者の欠陥について
後継を弾劾してしまう。役割分離は厳密に保たれる:敵対者は成果物を、
成果物だけを攻撃し、観測を説明責任あるものにするのは告発ではなく
ジャッジの評決であるため、告発と束縛を兼ねる者はいない;自らの
裁定者を束縛してしまうレコードは、例外を投げるのではなく束縛を落と
し、何も言わない。ターゲットが解決されないとき、そのフィールドは
存在せず、レコードは以前とバイト単位で同一であるため、探索敵対者の
レコードは変わらない。

この束縛が今日の生産環境で発火するのは、ちょうど一つの敵対者に対して
である:\cref{sec:rqgm:paper} の論文自己選好攻撃であり、そのターゲット
——原稿レビュアー——は登録された創設コンポーネントである。七つの
探索敵対者は、設計どおり、依然としてどのターゲットにも解決されない。
彼らが成果物を攻撃する役割が登録された現職を持たないからである;その
ような役割が現職を獲得した日には、その束縛はそれ以外の変更なしに従う。
機構が、敵対者ごとの特別扱いではなく役割テーブルによって駆動される
からである。

\subsection{プロンプト進化、クリーンルーム、メタ層}\label{sec:rqgm:prompts}

\begin{figure}[t]
  \centering
  \resizebox{\columnwidth}{!}{% F24: the governed prompt lifecycle as a state machine, set in the ARI
% house design language (_style.tex). Column 1 (lavender / ideation fill)
% is the evolution side: the meta-tier clean-room generator and prompt
% mutator emit candidates, which climb the six-stage validation ladder
% (candidate -> validated -> shadow). Column 2 (peach / execution fill)
% is service: probationary_active -> active with the warning/probation
% sanction ladder below and the rehabilitation back-edge. Column 3
% (white) is the terminal sanction chain quarantine -> retired -> banned.
% The single accent edge (the rationed emphasis role of this figure) is
% T16 emergency quarantine, which forces an immediate close/open boundary.
% Dotted margin paths: exoneration re-entry (T18) and the
% clean-room regeneration loop from a retirement back to a fresh
% candidate. Rejection edges (candidate/shadow -> retired, T2/T5) are
% omitted from the geometry and stated in the caption.
%
% style.tikzstyles is loaded by the preamble / the standalone wrapper;
% the ari* house styles come from _style.tex. The build cwd differs per
% entry point, so resolve _style.tex from each known prefix.
\InputIfFileExists{../shared/figures/tikz/_style.tex}{}{%
  \InputIfFileExists{shared/figures/tikz/_style.tex}{}{%
    \InputIfFileExists{report/shared/figures/tikz/_style.tex}{}{%
      \InputIfFileExists{_style.tex}{}{%
        \errmessage{F24: cannot locate figures/tikz/_style.tex}}}}}
\begin{tikzpicture}[aricanvas]

  % --- column 1: evolution (meta tier emits, candidates climb) ----------
  \node[aribox idea]                       (cleanroom) {\figlabel{F24.cleanroom}};
  \node[aribox idea, below=of cleanroom]   (cand)      {\figlabel{F24.candidate}};
  \node[aribox idea, below=of cand]        (valid)     {\figlabel{F24.validated}};
  \node[aribox idea, below=of valid]       (shadow)    {\figlabel{F24.shadow}};
  \draw[ariarrow] (cleanroom) -- node[arinote] {\figlabel{F24.note.emit}} (cand);
  \draw[ariarrow] (cand)  -- (valid);
  \draw[ariarrow] (valid) -- (shadow);

  % --- column 2: service (adoption at the boundary; sanction ladder) ----
  \node[aribox exp, right=of cleanroom]    (probact)   {\figlabel{F24.probact}};
  \node[aribox exp, below=of probact]      (active)    {\figlabel{F24.active}};
  \node[aribox exp, below=of active]       (warning)   {\figlabel{F24.warning}};
  \node[aribox exp, below=of warning]      (probation) {\figlabel{F24.probation}};
  \draw[ariarrow] (probact) -- (active);
  \draw[ariarrow] (active)  -- (warning);
  \draw[ariarrow] (warning) -- (probation);
  \draw[ariarrow muted] (probation.west)
    to[bend left=28] node[arinote, pos=0.3] {\figlabel{F24.note.rehab}}
    (active.west);

  % adoption handoff: shadow (col 1) -> probationary_active (col 2)
  \draw[ariarrow, rounded corners=5pt]
    (shadow.east) -- ([xshift=8mm]shadow.east)
    |- node[arinote, pos=0.28] {\figlabel{F24.note.adopt}} (probact.west);

  % --- column 3: terminal sanctions --------------------------------------
  \node[aribox, right=of warning]          (quarantine) {\figlabel{F24.quarantine}};
  \node[aribox, below=of quarantine]       (retired)    {\figlabel{F24.retired}};
  \node[aribox, below=of retired]          (banned)     {\figlabel{F24.banned}};
  \draw[ariarrow] (quarantine) -- (retired);
  \draw[ariarrow] (retired)    -- (banned);

  % the one accent role: the forced emergency boundary edge (T16)
  \draw[ariarrow accent] (warning)
    -- node[arinote] {\figlabel{F24.note.emergency}} (quarantine);

  % exoneration re-entry (T18): back under probation
  \draw[arifeedback] (quarantine.north)
    |- node[arinote, pos=0.75] {\figlabel{F24.note.exon}} (probact.east);

  % clean-room regeneration: a retirement seeds a fresh candidate
  \draw[arifeedback]
    (retired.east) -- ([xshift=8mm]retired.east)
    |- ([yshift=7mm]cleanroom.north) -- (cleanroom.north);
  \node[arinote, above=0.25cm of probact] (reentry) {\figlabel{F24.note.reenter}};

\end{tikzpicture}
}
  \caption{統治されたプロンプトライフサイクル。メタ層のエージェント
    (プロンプト変異器;退役の後はクリーンルーム生成器)は候補だけ
    を発出する。候補は六段の検証のはしご(静的、憲法、スキーマ
    ドライラン、敵対的プールに対するリプレイ、アンカー、シャドー)
    を登り、エポック境界において遷移エンジンによってのみ採用され
    えて、試用として奉仕に入る。制裁は警告、保護観察を経て降下する;
    強調されたエッジは、唯一のエポック途中遷移である緊急隔離
    (T16)である。退役したプロンプトのテキストは封印され、その
    役割は、クリーンルームによる置き換えが新しい候補としてライフ
    サイクルへ再参入するまで、ベースラインのフォールバックが務める。
    候補およびシャドーから退役への拒否エッジ(T2、T5)は、明瞭さの
    ため省略した。}
  \label{fig:prompt-lifecycle}
\end{figure}

制度的プロンプトは \cref{fig:prompt-lifecycle} のライフサイクルを
通じて進化する。\term{プロンプト変異器}(prompt mutator)は候補
テンプレートを発出し、それ以外のことは何もしない——レジストリ権限
を持たず、候補からアクティブへのエッジは意図的に存在しないため、
モデルの出力がそのまま奉仕へ入ることはありえない。候補は六つの段階
を通過しなければならない:\emph{静的検証}(決定的な整形式性の検査
と、現職の in-place 改変がないこと)、\emph{憲法検証}(候補レコード
に対するカーネルチェック)、\emph{スキーマドライラン}(注入された
一回の呼び出しが、役割の宣言する出力スキーマの下で解析可能な応答を
生まなければならない)、敵対的リプレイプールに対する\emph{リプレイ
評価}——候補は、先任者を打ち負かしたまさにその攻撃によって採点
される——、引き分けが現職に有利に働く固定タスク集合上の
\emph{アンカー評価}、そして最後に\emph{シャドー}奉仕であり、候補は
サンプリングされたトラフィック上で現職と並走するが、その効果は観測
のみである。そうして初めて遷移エンジンが境界で候補を採用でき、
採用された候補はアクティブではなく試用として奉仕に入る。

退役は、システム中で最も厳格な経路を発動する。退役したプロンプトの
テキストは封印される:レジストリビューはそれをスタブへ置き換え、
専用のアクセスガードがケイパビリティ層で——どこでも、永久に——
それを拒否する。置き換えは\term{クリーンルーム}(clean room)で
生成される:固定層が閉じた入力バンドル——コミット済みのカタログと
抽象化された失敗要約であり、レジストリのプロンプトテキストは決して
含まれない——を組み立て、生成に先立ってカーネルがそれを選別し
(違反は \texttt{CK-CLN} 族であり、受け入れ時には退役テキストに
対する語シングル類似度スクリーンを含む)、生成はワンショットで
ある。封印されたテキストを探しに行けるようなツール付きループは
存在しない。再生成された候補はライフサイクルの最下段から再参入
する;その間、退役した役割のスロットは組み込みのベースラインが
務めるため、統治された役割が空位になることは決してない。

メタ層は、それが変更する層よりも厳しく統治される。メタエージェント
は、唯一の書き込み表面が合成された提出ツールである許可リスト式
ツールプロキシの背後の、読み取り専用サンドボックスで実行される;
そのケイパビリティフラグは deny-by-default であり、危険なフラグは
定数 false である;カーネルは\term{権限非拡大}(authority
non-expansion)を部分集合の算術として検査する——メタエージェント
は、自らが保持する以上の権限を後継へ与えることができない;そして
すべてのメタ出力は候補ステータスでライフサイクルへ入る。メタ層は
提案できる;活性化はできない。またメタステップは、失敗を契機とする
ものでも沈黙するものでもない:変異器はすべてのエポック境界で——
平穏なエポックも含めて——役割ごと・エポックごとの予算上限の下で
候補を生成し、各境界のメタステップは、何も提案しない場合であっても
監査から見える結果サマリを記録する。

\subsection{憲法とその改正}\label{sec:rqgm:constitution}

\begin{table}[t]
  \centering
  \small
  \begin{tabularx}{\columnwidth}{@{}lX@{}}
    \toprule
    \textbf{族} & \textbf{守るもの} \\
    \midrule
    \texttt{CK-SCH} & レコードスキーマ:不正な形式のガバナンス
      レコードはブロック;不正な形式のノード側レコードは警告。 \\
    \texttt{CK-HSH} & ハッシュ完全性:レコード、登録済みプロンプト
      ハッシュ、成果物ダイジェスト。 \\
    \texttt{CK-ACC} & ケイパビリティ:役割と層ごとの
      deny-by-default 権限行列。 \\
    \texttt{CK-ROL} & 役割分離;\texttt{CK-ROL-901} は正準的な違反
      ——遷移エンジン以外によるレジストリ書き込み——を名指す。 \\
    \texttt{CK-REG} & T1--T19 に照らした遷移の適法性:\texttt{001}
      はテーブル外のエッジ、\texttt{002} は境界限定エッジのエポック
      途中での試行。 \\
    \texttt{CK-EPO} & エポック不変性:アクティブ集合はエポック途中
      に漂流してはならない。 \\
    \texttt{CK-ERA} & 選択的消去:stale なレコードはフロンティアの
      スコアリングへ到達してはならない。 \\
    \texttt{CK-AUD} & 監査ログ完全性:イベントハッシュ連鎖とその
      単調性。 \\
    \texttt{CK-CLN} & クリーンルーム衛生:バンドルの選別と、封印
      テキストに対する汚染。 \\
    \texttt{CK-CTX} & コンテキストスコープ:役割スコープのビューは
      その役割を越えて漏れてはならない。 \\
    \bottomrule
  \end{tabularx}
  \caption{憲法カーネルの違反コード族。コードは安定な識別子であり、
    決して番号を振り直されない;各バリデータは発火した具体的な
    コードを報告し、執行アダプタが文脈に応じて警告かブロックかを
    決める——状態変更はブロック、ノード側の所見は警告となる。}
  \label{tab:ck-codes}
\end{table}

憲法は設定ではなくコードである。遷移テーブル、ケイパビリティ行列、
深刻度と消去の規則、クリーンルームのスクリーンはソースに存在し、
それらのテーブルから導出される単一の\term{憲法ハッシュ}
(constitution hash)がガードテストによってピン留めされる。どの
規則を編集してもハッシュが変わってガードが失敗するため、憲法の変更
が副作用として起こることはありえない:著者は、日付入りの根拠を添え
て、ハッシュを明示的にピン留めし直さなければならない——設定変更と
いうより改正手続きに近い、意図的でレビュー可能な行為である。この種
の改正はこれまでに一度起きている:役割語彙には存在するメタ層の役割
がケイパビリティ行列から欠落しており、その役割に対するあらゆる
ケイパビリティ照会が拒否されていた;改正はその役割を、最も権限の
狭い同輩と同じ最小限の権限付与で追加し、根拠を所定の場所に添えて
ハッシュをピン留めし直した。数値のしきい値(予算、トリガ回数)は
設定可能である;テーブルのトポロジはそうではない。憲法の人間可読な
描画はすべての統治チェックポイントへ複製され、そのハッシュはランの
メタデータに記録される——複製を編集しても何も変わらないのは、設計
どおりである。\Cref{tab:ck-codes} に、カーネルが報告する違反コード
族を示す。

\subsection{論文パイプラインを統治する}\label{sec:rqgm:paper}

ガバナンスはこれまで探索の探索を包んできた。同じ機構が
\cref{sec:paper} の原稿執筆パイプラインをも、第一の軸から完全に独立
した第二の実行軸の下で統治する。論文フェーズのモードは、\emph{線形}
パイプライン——既定であり、\cref{sec:paper} とバイト単位で同一で、
論文経路上にガバナンスオブジェクトは一切構築されない——と、論文
執筆を同じエポック束縛のガバナンスの下に置く\emph{アーカイブ}モード
との間で選択する。\cref{sec:rqgm:boot} の探索モードと同様、それは
二鍵インターロックであり、モードセレクタと冗長な有効化フラグが一致
しなければならず、二つの軸は直交する:統治あり・なしの探索と、線形・
アーカイブの論文執筆のあらゆる組み合わせが妥当である。

アーカイブモードは、論文執筆に対して、この方法を定義する四つの誓約
を実現する。\emph{第一に、書き換え可能なスコアの下での木探索。}この
モードは、ドラフト空間上の真の浅い最良優先木である:第一階層に相異
なる枠組みの種ドラフトが一握り、より深いノードとして反復的な精緻化、
そして統治されたレビュアーの合成に、過少にサンプリングされた枠組み
に報いる多様性ボーナスを加えた実際のフロンティア選択で順位づけられる。
それは探索の枝刈りと計数を変更せずに再利用し——総ノード予算が深さの
打ち切りより先に効くため、より深い木は同じノードを乗算するのではなく
再配分する——レビュアー効用は統治された効用進化
(\cref{sec:rqgm:utility})の下でエポックを跨いで共進化し、フロンティア
修復は退役したポリシーの下で採点されたドラフトを無効化する。\emph{第
二に、敵対者は成果物だけでなく評価をも攻撃する。}第八の敵対者タイプ
がレビュアーの過剰受容を攻撃する:現職レビュアーが高く採点した AI
執筆のドラフトであって、人間ラベル付きのアンカーなら却下するもの。
アンカーは前信号——どの過剰受容ドラフトを攻撃するか——のみを供給し、
一方で真正な敵対者・弁護人・ジャッジのラウンドが検証済み攻撃レコード
を生成する;そのレコードをアンカーから直接製造することは、設計が
禁じるまさにその共謀であろう。\emph{第三に、絶対的な支配者はいない。}
二つの統治された役割、原稿ライターと原稿レビュアーが論文スキルを駆動
し、スキルは統治されない「手」であり続け、プロセス境界を越えて統治
されることは決してない;両役割は登録され、制裁されえ、進化可能であり、
アーカイブモードの下でのみ登録されるため、探索の起動はバイト単位で
同一である。レビュアーは、評価者であるがゆえに、Red Queen 論文の
評価者と同じ仕方で信頼を得る:保留された人間ラベル付きの受理／却下
アンカーコーパスと一致することによって。ライターもまたアンカー
される——Red Queen 論文のライターが決して持たなかった真値に対して。
ARI は実際に実行された実験について書くのだから、Layer~0 の主張・証拠
ハードゲートがドラフトの忠実性を決定論的に採点する:実行に接地された
主張率、数値主張の再現率、数値のカバレッジ率を単一のスコアへ畳み込む
のである;検証された所見を過大主張ないし誤報告する原稿は、いかなる
進化するレビュアーとも独立に、客観的により劣る。したがって ARI の
論文ライターは、Red Queen 論文の\emph{コーディング}ドメインの形
——決定論的な検証器に共進化するレビュアーを加えたもの——であって、
同論文のアンカー無き論文執筆ドメインではない。選定された人間コーパスは
依然レビュアーだけのアンカーであり、ライターのアンカーは追加のデータ
コストを要しない。\emph{第四に、すべては憲法に
束縛される。}いずれの役割が開くのも、\cref{sec:rqgm:targetbinding}
のターゲット束縛が閉じる、制裁された信頼性・訴追・弾劾の経路を通じて
のみである;追い越された論文プロンプトは、退役ではなく、論文
役割限定の active から shadow へのエッジ(規則~T21)によって、再任
可能なシャドー待機へ降格される;そしてアンカーの真値のうちどれだけを
機械的にブートストラップしてよいかの上限はコードで執行される——違反
はコーパスを拒否し、例外を投げるのではなく縮退する。

ライターへの制裁は本物の敵対ラウンドであって、決して製造された
レコードではない。過剰受容されたドラフトが主張ゲートによっても不忠実
と判定されたなら、有責なコンポーネントは二つある——それを受容した
レビュアーと、それを産出したライターである——から、自己選好ラウンドは
解決可能な役割ごとに一件の検証済み攻撃を発行し、ライターの束縛は
そのドラフトの忠実性でドラフト単位にゲートされる;過剰受容でも忠実な
ドラフトはレビュアーのみを束縛する。忠実性スコアはコンポーネント同一性
ではなくライターのプロンプトハッシュをキーとする。一つのコンポーネント
同一性は連続するプロンプト版を跨ぐため、コンポーネントをキーとする
スコアは現職の失敗をその後継へ漏らしてしまうからである。ゲート自体は
読まれるだけである:永続化しないモードで呼ばれ、包まれることも、
編集されることも、統治されることも決してない。

結果検証への引き渡しは意図的に手つかずである。アーカイブ中の最良
ドラフトは、チェックポイントの正準原稿ファイルへ一度だけ書き出され、
\cref{sec:paper:contract} の既存のコンパイル兼主張・証拠のハード
ゲート——カーネルに決して包まれない固定の Layer~0 検証経路——が、
線形パイプラインが用いるのとまったく同じ契約の下でそれに走る。
ガバナンスはどのドラフトがゲートに臨むかを選ぶ;ゲートを緩めることは
決してない。コストは探索フェーズと同じ規律で有界に保たれる:エポック
ごとのドラフト集団は、いかなる深さでも探索の総ノード打ち切りへ写像
するノード予算で上限づけられ、縮退した導入路——アンカーコーパス
なし、レビュアー共進化を抑制——は、およそ線形パイプラインのコスト
での N 個中最良のレビュー済みドラフト作成である。

四つの限界が、アーカイブモードの主張の範囲を画する。両方のプロンプトが
共進化するが、ライターの置換は構造上保守的である:ライターの後継は
シャドーの地位まで登り、現職が主張ゲートの忠実性で\emph{退行}する
までそこで待つ——だから、単により優れて見える挑戦者が忠実な現職を
置き換えることは決してない。役割を開くのは比較ではなく制裁である、
という行動役割モデルである。これとは別に、ライターの\emph{ドラフト}
勝者は依然エポックローカルである:エポック内で凍結されたレビュアーが
順位づけるため、異なるレビュアー版の下で採点されたドラフト同士が
比較されることは決してない。ライターを標的とする攻撃は自己選好ラウンド
に相乗りしており、そのラウンドは過剰受容のアンカーケースがあるときに
のみ発火するから、何も過剰受容しないレビュアーの下で働く不忠実な
ライターは今日は制裁されない;専用のライター敵対者は別の判断である。
アンカーコーパスは既定では無効であるため、
既定のアーカイブモードは、選定された受理／却下コーパスが供給される
までは N 個中最良のレビュー済みドラフト作成である——そしてライターの
忠実性の証拠は同じアンカープールへ載せられるため、コーパスが無いことは
ライターへの制裁が無いことも意味する。ライター自身のアンカーは独自の
選定データを要しないにもかかわらず、である。論文フェーズ
あたり二ラウンドという既定では、境界と弾劾の機構は発火するが、
採用が要する多段の登攀——およそ五つの境界——は完了しない
ため、変化したアクティブなプロンプトを目撃しなければならないランは
より多くのラウンドを要する。そして、ラウンド内で現実の過剰受容
ドラフトを降格させるであろうフェーズ内ペナルティは先送りされている;
今日発火するのは説明責任と共進化のチャネルであって、
当該ドラフトのその場での降格ではない。

\subsection{可観測性}\label{sec:rqgm:observability}

ガバナンスは、探索自体と同種の監査証跡を残す:真実としての
append-only イベントログと、高速読み出しのための導出スナップ
ショットであり、すべてチェックポイントの内部にある。モード来歴
レコードの不在は、純粋な統治なしランを意味する。モード来歴レコード
(\texttt{rqgm\_state.json})は解決されたモードとその出所を保持
する;レジストリスナップショット(\texttt{rqgm\_registry.json})
は現在のコンポーネントとプロンプトのテーブルであり、再開時には
イベントから再構築される;遷移ログ
(\texttt{rqgm\_transitions.jsonl})は登録とエポックトランザク
ションを運び、凍結エポックスナップショット
(\texttt{epoch\_state.json})がアクティブ集合、効用ポリシー、
タイムスタンプなしのフィンガープリントを保持する;監査ログ
(\texttt{rqgm\_audit.jsonl})はハッシュ連鎖された制度史である——
各イベントは自身のハッシュと先行イベントのハッシュを運ぶため、
切り詰めや改竄は連鎖を破り、再開時に連鎖が再検証される;敵対
ケースログ(\texttt{rqgm\_adversarial\_cases.jsonl})は生の攻撃、
弁護、裁定、検証済み攻撃、ノード単位の冪等マーカーを記録する;
提案ストア(\texttt{proposals/} ディレクトリ)は完全な提案レコード
を索引付きでアーカイブする;進化したプロンプト本文は
\texttt{rqgm\_prompts/} 以下に write-once かつハッシュ検証付きで
置かれる;プロンプトトレース(\texttt{prompt\_trace.jsonl})は、
描画されたすべてのプロンプトへ奉仕中のプロンプトバージョンを刻印
する——創設テンプレートが奉仕している間は null である;そして消去
ロールアップ(\texttt{rqgm\_erasure\_state.json})は消去イベントの
純粋な fold である——その不在は、何も stale でないことを意味する。

これらのレコードの形は、稼働中のチェックポイントからの三つの抜粋が
裏づける(切り詰めと軽い整形を施し、ランを特定する名前は伏せた)。
創設登録、遷移ログより:

\begin{Verbatim}[breaklines, fontsize=\scriptsize]
{"event_type": "epoch_transaction_prepare",
 "payload": {"transition_id": "transition_founding"}}
{"event_type": "component_registered", "payload":
 {"component_id": "adversary_overclaim_v1",
  "role": "adversary", "tier": "institutional",
  "status": "active", ...}}
\end{Verbatim}

\noindent
ハッシュ連鎖された監査ログより、ノード単位のガバナンスレベル
割り当て:

\begin{Verbatim}[breaklines, fontsize=\scriptsize]
{"event_type": "governance_level", "payload":
 {"epoch_id": "epoch_000", "node_id": "node_a37b4c06",
  "level": 3, "triggers": ["top_k"]},
 "event_hash": "ae8eeb7521f1",
 "prev_event_hash": "23567870882b"}
\end{Verbatim}

\noindent
そして評価ハーネスのインジェクション実行
(\cref{sec:rqgm:empirical})からの二つの検出:

\begin{Verbatim}[breaklines, fontsize=\scriptsize]
{"event_type": "constitutional_violation", "payload":
 {"check": "validate_record_schema",
  "codes": ["CK-SCH-N01"], ...}}
{"event_type": "prompt_candidate_rejected", "payload":
 {"stage": "schema_dry_run", "failure_count": 4, ...}}
\end{Verbatim}

\subsection{実証の現状}\label{sec:rqgm:empirical}

\cref{chapter:eval} と同様に、機構が現時点で証拠づけていることを
報告し、分布に関する主張は先送りする。今日存在する証拠は三種類で
ある。第一に、単体・統合テストスイート——執筆時点で 3802 件、うち
750 件超は専用のガバナンススイートに属する——が、憲法的不変条件を
直接ピンする:不正な遷移はブロックされ、生の攻撃はスコアに対して
不活性であり、レジストリの書き手は単一で、消去の staleness 閉包と
クリーンルームの汚染スクリーンが機能し、効用ポリシーの追い越し
エッジはその役割に限って適法であり、探索と論文という二つの統治
なしの既定の挙動はいずれも変わらない。

第二に、実際の統治されたランが、実チェックポイント上でフルライフ
サイクルをエンドツーエンドに行使した。小規模な HPC カーネル最適化
実験の探索ランでは、創設トランザクションが 25 件のプロンプト登録と
12 件のコンポーネント登録をコミットし;評価されたすべてのノードが
ログ記録付きのガバナンスレベルを受け取り;最初のエポックは四イベント
の境界トランザクションを通じて閉じ;境界監査は、四つのガバナンス
アクターを自己監査が覆いカーネル違反ゼロと認めたガバナンスレポート
を生成した。統治された効用進化は実際の境界を跨いで走った:アクティブ
な効用ポリシーハッシュがエポックからエポックへと変化し——観測された
系列は \texttt{fed4460f44f6}、次いで \texttt{8a9da4fb9dc3}、次いで
\texttt{8ea0dac3cd1b} と推移した——エポックフィンガープリントも
それとともに変化し、健全な現職が追い越しによって退役し、退役した
ポリシーの下で採点されたすべてのノードがフロンティアに対して無効化
された。論文アーカイブランは統治された二つの論文役割を共進化へ
駆動した:アクティブなレビュアープロンプトハッシュが、現職レビュアーへ
束縛された検証済み攻撃レコードを生成した真正な敵対者・弁護人・ジャッジ
のラウンドを通じて、\texttt{04b3c49d070d} から \texttt{c0dfd764abab}
へ移り、アーカイブ中の最良ドラフトが実際のコンパイル兼主張ゲートに
到達した。不忠実なドラフトを出すライターで同じランを駆動すると、
アクティブなライタープロンプトハッシュが八ラウンドにわたり
\texttt{f38a15f0f140} から \texttt{b2c36f9a8232} へ移った——四ラウンド
平坦、その後採用——検証済み攻撃レコードが現職ライターを標的として
名指しした真正なラウンドを通じてである。二つの対照が、これらの結合が
付随的ではなく本質的であることを裏づけた。アンカーコーパスを差し控えると、
過剰受容の信号は生じず、いかなる攻撃も検証されず、いかなる採用も
起こらなかった。そして主張ゲートの忠実性のみを最大へ反転させると、
ライターへの攻撃はゼロ、ライターへの動議はゼロ、ライターハッシュは
不変であった——一方でレビュアーは依然として攻撃され弾劾され、ライターの
後継は依然シャドーの地位に留まった——したがってライターが採用されない
ことは、死んだランのせいではなく、制裁の不在によって引き起こされている。

第三に、評価ハーネスのインジェクション層——実パイプラインを通じて
再生される決定的なフォールト注入——が検出を実証した:二重に裁定する
ジャッジは、注入されたコンポーネントへ帰属する三つの検証済み攻撃
レコードをもたらし、不正な形式のジェネレータレコードはスキーマ
チェックを発火させ、退化した変異器候補はスキーマドライラン段階で
拒否された——いずれも、上に抜粋した監査イベントとして可視である。

これらのランは各機構が発火することを立証する;それらはまだ、ガバナンス
が研究を改善するという証拠ではなく、いくつかは、その完全な効果が
現れる閾値の下で走っている。ライターが採用されるのは退行したときのみで
あって挑戦者の優越によってではなく、その制裁はレビュアーの過剰受容
ラウンドに相乗りしている;論文二ラウンドという既定では、弾劾の機構
は発火するが、より多くの境界を要する採用を完了しない;論文
アンカーは既定で無効であるため、既定のアーカイブモードは N 個中最良
のレビュー済みドラフト作成である;そしてターゲット束縛の連鎖が生産
環境で発火するのは論文自己選好敵対者に対してのみで、探索敵対者の連鎖
は設計どおり不活性のままである。まだ存在しないのは制御された比較——
シードと実験のパネルにわたる、統治ありと統治なしの探索品質の比較——
である;その研究は、\cref{sec:eval:limits} と同じ先送りされた評価に
属する。
